\color{unidarkgreen}{\subsection{Настройка каналов напряжения}}
\color{black}

\begin{enumerate}

\item
Для подключения вторичных цепей напряжения от измерительных трансформаторов напряжения в устройстве предусмотрены следующие аналоговые входы:
\begin{itemize}
    \item 3 входа для подключения цепей напряжения от основной вторичной обмотки ТН (Ua, Ub, Uc);
    \item 1 вход для подключения цепи отбора напряжения (Uотб).
\end{itemize}

\item
Рабочий ТН, трехфазная система которого используется для работы защит и автоматики может быть
установлен как на линии, так и на шинах. Для корректной работы логики переключатель \textbf{«Место ТН»} должен
быть установлен в соответствующее положение.

\item
Аналоговый вход отбора напряжения, на который может быть подано напряжение от однофазного ТН
или ШОН, предназначен для логики АПВ.

\item
Схема настройки аналоговых значений Uэ1, Uэ2 в зависимости от схемы РУ приведена на рисунке \ref{pict_nkn:nkn}. Параметры настройки каналов напряжения представлены в таблице \ref{app_set:schema}.

\vspace{2mm}
\begin{figure}[H]
\centering
\includegraphics[width=1\textwidth,height=1\textheight,keepaspectratio]{\fbpath/10.01 АУВ/НКН/img/nkn.pdf}
\captionof{figure}{Схема настройки аналоговых значений Uлэп, Uэ1, Uэ2}\label{pict_nkn:nkn}
\end{figure}

\item
Если, например, рабочий ТН установлен на шинах, то значению Uэ2 присваивается одно из фазных либо
линейных напряжений, выбираемое переключателем \textbf{«Фаза подключения Uэ2»}. Переключатель \textbf{«Напр. Uэ2»}
в таком случае должен быть в положении \textit{«Фаза подкл.»}. А значению UЭ1 присваивается напряжение,
поступающее с аналогового входа отбора напряжения Uотб, скорректированное по амплитуде и фазе
соответствующими уставками \textbf{«К подстройки Uотб»} и \textbf{«Ф подстройки Uотб»} по следующей формуле:
\begin{align}
\vec{U}_{\text{отб корр}} = \vec{U}_{\text{отб}} \cdot K \cdot e^{j\phi},
 \end{align}
\begin{eqexpl}
  \item{$ K $} коэффициент подстройки Uотб по амплитуде;
	\vspace{2pt}
  \item{$ \phi $} угол подстройки Uотб по фазе.
\end{eqexpl}

\item
Если рабочий ТН установлен на ЛЭП, то наоборот: фаза подключения выбирается для Uэ1, а значению
Uэ2 присваивается скорректированное напряжение с аналогового входа отбора напряжения. 

\end{enumerate}